% Generated draft: requires mathematical and source review.
\documentclass[11pt]{article}
\usepackage[margin=25mm]{geometry}
\usepackage{fontspec}
\setmainfont{Latin Modern Roman}
\usepackage{amsmath,amssymb,mathtools,unicode-math}
\setmathfont{Latin Modern Math}
\usepackage{enumitem,needspace,xurl,hyperref}
\setlength{\parindent}{0pt}
\setlength{\parskip}{5pt}
\setlength{\emergencystretch}{3em}
\newcommand{\N}{\mathbb{N}}
\newcommand{\Z}{\mathbb{Z}}
\newcommand{\Q}{\mathbb{Q}}
\newcommand{\R}{\mathbb{R}}
\newcommand{\C}{\mathbb{C}}
\newcommand{\F}{\mathbb{F}}
\newcommand{\A}{\mathbb{A}}
\newcommand{\PP}{\mathbb{P}}
\newcommand{\E}{\mathbb{E}}
\newcommand{\eps}{\varepsilon}
\newcommand{\dd}{\,\mathrm{d}}
\begin{document}
\section*{UnsolvedMath ID: 30004501}
\textbf{Source problem code:} catalogue entry
\subsection{Definitions and mathematical statement}
Put \(k=\F_q(t)\), where \(q\) is a fixed prime power, and fix a quaternion algebra \(D\) over \(k\). For an extension \(L/k\), write \(D_L=D\otimes_k L\), let \(\A_L\) denote the adele ring of \(L\), and interpret \(D^\times(\A_L)\) as \((D_L\otimes_L\A_L)^\times\). Write \(Z\) for the center of this multiplicative algebraic group. Scalar extension gives the inclusion used to evaluate functions on \(D^\times(\A_k)\).

Let \(F/k\) range over quadratic extensions of genus \(g\), and let \(K\subset D^\times(\A_F)\) be a compact subgroup. Set
\[
H_K=K\cap D^\times(\A_k).
\]
The supplied statement describes \(m\) as the index of \(H_K\) in a maximal compact subgroup of \(D^\times(\A_k)\): thus, whenever a containing maximal compact subgroup \(M\) and a finite index are specified, the intended notation is
\[
m=[M:H_K].
\]
Neither the choice of \(M\) nor a hypothesis ensuring that this index is finite is supplied.

Let \(f\) be a normalized Hecke eigenfunction on
\[
X_{F,K}=D^\times(F)\backslash D^\times(\A_F)
       /\bigl(KZ(\A_F)\bigr).
\]
Here the meaning of ``normalized,'' the relevant Hecke operators, and the admissible automorphic function class are not specified in the supplied record.

\textbf{Proposed statement (with unresolved conventions).}
Do there exist constants \(C,\delta>0\), depending only on \(q\) and \(D\), such that, for every such \(F\), \(K\), and normalized Hecke eigenfunction \(f\), and every \(y\in D^\times(\A_k)\),
\[
f(y)=
\int_{D^\times(k)\backslash D^\times(\A_k)}
       f(x)\,\dd\mu(x)
+O\!\left(m^Cq^{(1-\delta)g/2}\right),
\]
with the implicit constant independent of \(F\) and \(K\)? The occurrences of \(f(x)\) and \(f(y)\) mean evaluation after scalar extension and projection to \(X_{F,K}\).

The measure \(\mu\), its normalization, and convergence of the displayed period are unresolved. In particular, \(f\) is invariant under the adelic center, whereas the displayed integration domain has no central quotient; a convention addressing central volume or convergence is needed. The record also does not explicitly settle whether the implicit constant is uniform in \(y\) and in the choice of normalized \(f\). No additional hypotheses or conventions are imposed here to resolve these issues. The neighboring source record cited as providing the common setup was not included in the supplied evidence.
\subsection{Short English statement}
For quadratic function fields \(F/\F_q(t)\), the question asks whether every normalized Hecke eigenfunction on the indicated quaternionic automorphic quotient differs at every base-field adelic point from its base-field period by \(O(m^Cq^{(1-\delta)g/2})\), with \(C,\delta>0\) depending only on \(q,D\) and the implicit constant independent of \(F,K\). The normalization, finite-index condition, period measure and convergence, and full uniformity conventions require clarification. This is a statement-only draft requiring human mathematical review.
\subsection{Sources}
\begin{itemize}
\item Local UnsolvedMath registry, supplied source record: permanent ID \(30004501\), exact source problem code \texttt{catalogue entry}, title ``Pointwise Function-Field Equidistribution from Period Integrals,'' in the Oberwolfach Reports Open Problems collection. No individual proposer is supplied; the proposed year is recorded as 2020.
\item Supplied report citation: \emph{Automorphic Forms and Arithmetic} (2021), Oberwolfach Reports, DOI \texttt{10.4171/owr/2020/26}, pp.\ 1317--1383. Supplied source and literature link: \url{https://doi.org/10.4171/owr/2020/26}. The supplied assessment identifies the question as Question 6 in Oberwolfach Report 26/2020. The background describes these reports as records of mathematical workshops and their open-problem discussions.
\item Supplied literature assessment, checked on \texttt{2026-08-22}: the universal function-field asymptotic with a power saving is described as open; the report is said to give only a split-case calculation and a suspicion that the formula is correct. This is supplied dated curation, not an independently verified literature assessment. The supplied background explicitly requests independent verification of time-sensitive status and cited claims.
\item The registry labels the statement \texttt{corrected\_verified} and records that the official report was checked, that the common setup comes from the canonical neighboring record \texttt{OWR-1703872-006}, and that the exact dependence of all constants was retained. These are supplied verification claims; neither that neighboring record nor the report text is included here, and no independent source verification is asserted.
\end{itemize}
\end{document}
